\maketitle
%\thispagestyle{fancy}
\section*{Overview}

\begin{enumerate}
  \item \emph{Price stability.}
    Created options and budget estimates for price stability work.

    Argued in favour of changing the price oracle directly.  Created a notebook that simulates the effects of some simple alternative protocol rules on the past 12 months of storage pricing: \url{https://dusky-harbor-6a1f04.swarm.shtuka.io/}

  \item \emph{Protocol risk.} 
    Began design work on protocol risk framework.

\end{enumerate}

\section*{Project management}

\subsection*{Completed}

\begin{itemize}
  \item Pricing: First exploration of price controller design space. Present options.
\end{itemize}

\subsection*{Ongoing}

\begin{itemize}
  \item Upgrade process: streamlined stake migration and redistribution upgrade flow.
  \item Protocol risk: Develop risk model for supply-side stability under protocol changes.
  \item Hardening consensus: analyse weaknesses, evaluate threat level, list mitigations.
  \item Pricing: Build alignment around user-facing price objectives framework. 
\end{itemize}


In our design space exploration, we identified solutions that address price stability but also address objectives about price \emph{level} that we had not thus far considered.
%
Consequently, we are reframing the ``price stability'' workstream to be more generally about storage pricing.

\subsection*{Forthcoming}

\begin{itemize}
  \item Pricing: Build alignment around user-facing price objectives framework. 
\end{itemize}

\subsection*{Backlog}
\begin{itemize}
  \item Upgrade process: revisit migration-free options for stake registry.
  \item Negative incentives: explore and evaluate solutions.
  \item Node dispersion: build mathematical model for address allocation schemes and node dispersion.
  \item Big nodes: Establish framework for converting findings into contributions back into mainline bee.
\end{itemize}


\section*{Current work}

\subsection*{Swarm Network upgrade process}

We've been talking about rolling out changes to Swarm's chain protocol for some time. 
%
It's time we revisited the question: \emph{How do we actually execute these upgrades?}

Swarm's state is maintained across a decentralised network of nodes, which makes upgrading no easy feat.
%
We proposed the following recipe:

\begin{enumerate}
  \item \emph{Redeploy contracts every time.} Every time we make a breaking change to the protocol, whether wire-only or onchain, redeploy the full contract suite.

  \emph{Why?} Incentivising operators to switch to the new version means we have to stop paying them to run the old version and start paying them to run the new version.
  %
  A single \code{Redistribution} contract cannot tell which commits come from the new version, so in order to redirect revenue we need to deploy a new one.
  %
  Partial redeployments require an admin privilege with a large attack surface that we should aim to renounce; if we do this, redeploying \code{Redistribution} means redeploying all incentive contracts.

  \item \emph{Declare a cutover block height.} Announce the target block height at which the network will cut over to the new protocol version.

  \emph{Why?} We need to provide operators with a clear social contract: upgrade within the specified window, and you will receive no interruptions to your income.
  %
  This contract must be communicated clearly in advance, with plenty of time left to react.

  \item \emph{Put the cutover in the client.} Integrate the cutover into clients: the latest bee version should know how to speak both current and upcoming protocol versions, and how to migrate stake from old to new.
  
  \emph{Why?} In order to offer a consistent contract to operators upgrading within a window, client behaviour needs to be the same regardless of when in that window the upgrade happened.
  %
  Old clients and upgraded clients must have the same behaviour within the window.
  %
  That means letting the client handle when to switch protocol versions rather than making it an operator decision.

  \item \emph{Make the cutover atomic.} Broadcast a single transaction  that pauses the old contracts and unpauses the new contracts at the same instant. As soon as clients observe this transaction, they execute the cutover.

  \emph{Why?} Payments to operators on the old version must stop at the same time we start paying operators on the new version.
  %
  If we stop it later, so that the incentives for two versions overlap, we get perverse incentives and piecemeal migration.
  %
  If we stop it earlier, we get downtime during which no operator gets paid.
  %
  Given that we're doing it at the same time, we can derisk the operation by making it atomic.
\end{enumerate}


\begin{figure}[ht]
  \begin{center}
  \includegraphics[width=\linewidth]{coordinated-fork-migration}
  \caption{Coordinated migration across a protocol cutover. State is represented as a bitmap, which we can think of as a toy model of the storage index. Yellow circles are bees that haven't yet upgraded; green circles are bees that have.}
  \label{fig:coordinated-fork-migration}
  \end{center}
\end{figure}


\paragraph{Transition challenges}
Redeploying all contracts entails two challenges that we haven't always had to deal with: migrating stake and migrating batches.
%
If we're going to do this on every upgrade, we need to make this operation seamless and low-risk.
%
For a step in this direction, we proposed integrating stake migration into the client: \url{https://github.com/ethersphere/bee/issues/5577}.

While we work towards a solution for migrating batches, we can carry out surgical redeployments of \code{Redistribution} and any other contracts whose code has changed, authorising the new deployment to redistribute revenue from the existing \code{Postage} contract at the same time as pausing the old ones.
%
The same arguments for making the cutover atomic and putting it in the client apply.

This proposal doesn't address the question of how to transition the wire protocol at all. That is another story for another day.

\paragraph{Alternate proposals}
%
An alternative approach that has been proposed is to recentralise upgrades by putting all contracts behind an upgradable proxy: \url{https://github.com/ethersphere/storage-incentives/pull/310}. 

We argue that putting stake behind a proxy increases the attack surface for admin keys (from ``burn all stake'' to ``steal all stake'').
%
It's also not clear that there is a pressing need to proxy stake, pricing, or redistribution; of these, only stake needs significant state and assets to be transferred, and stake migration UX can be improved relatively easily with changes at the client level.
%
We support the idea of putting batches behind a proxy until the means to coordinate a user-driven batch migration can be established.